\section{Upper-bound classification: no seventh irreducible role}\label{sec:upper-bound-classification}

% Mathematical job:
%   - Search candidate seventh-role features in the covered R010 scope.
%   - Apply the repaired role-projection alignment rules to each candidate.
%   - Close the three high-priority threats:
%       probability / uncertainty,
%       residual causality,
%       residual agency.
%   - State the upper-bound classification theorem within the covered scope:
%     every candidate feature classifies as one of
%       {P1..P6 projection, composite, presentation variant, level shift,
%        activation-threshold refinement, instrument-visibility artifact,
%        empirical-bridge annotation, outside-domain structure,
%        already-covered bookkeeping/annotation}.
%   - No covered candidate remains as an unresolved seventh-role threat or
%     as a genuine seventh irreducible typed closure role.
%
% Role-alignment non-negotiables (carry forward from Section 2):
%   NOT P3 for generic transition.
%   NOT P5 for generic path / derivation.
%   NOT P1 for generic semantics / interpretation.
%
% Governing sources:
%   workspace/meta/lens_lab/campaigns/paper3_exact_six_exhaustiveness/upper_bound_classification_repair/T304_content_review.md
%   workspace/meta/lens_lab/campaigns/paper3_exact_six_exhaustiveness/upper_bound_classification_repair/role_projection_alignment_rules.md
%   workspace/meta/lens_lab/aristotle/results/paper3/paper3_upper_bound_classification_repair_F/paper3_upper_bound_classification_repair_F_result_shape.md
%   docs/papers/Tsiokos_2026_Six_Birds_Protocol_Trap_Holonomy_Without_Entropy_Production.tex
%     (P3-is-a-diagnostic principle)
%   paper/prior_literature_context.md (Tier A3, B3)
%
% Overclaim risks:
%   - Claiming no-seventh-role OUTSIDE the covered scope.
%   - Dismissing a candidate by definition rather than by classification.
%   - Backsliding into the forbidden role assignments above.
%   - Treating classification pass as proof of uniqueness.


The lower bounds of Section~\ref{sec:lower-bounds} rule out the
trivialization in which some primitive role lacks an admissible
witness in \FATCD{}. The upper-bound direction addresses the
complementary question: whether some additional, seventh closure role
is needed in the covered scope. This section shows that no such role
is needed. The argument proceeds in three steps. First, we state the
alignment rules that constrain how raw features may project to
\Pone{}, \ldots, \Psix{}. Second, we close the three high-priority
candidate seventh-role threats: probability and uncertainty, residual
causality, and residual agency. Third, we record the upper-bound
classification theorem that collects these local classifications into
a single covered-scope statement.

\subsection{Role-projection alignment rules}\label{subsec:alignment-rules}

The role meanings of Definition~\ref{def:primitive-roles} carry
alignment rules constraining which raw features may be interpreted as
active projections of each role. These rules refine the corrected
alignments of Remark~\ref{rem:repaired-alignments} and form the
discipline under which the classification theorem is proved.

\begin{proposition}[Role-projection alignment]\label{prop:alignment-rules}
Under honest bookkeeping and the witness discipline of
Definition~\ref{def:common-witness-data}, each role projection is
constrained as follows.
\begin{itemize}[topsep=2pt,itemsep=1pt]
  \item \Pone{}: admissible only when the candidate cluster carries a
    selected update, a quotient or relation record, and a witness of
    descent compatibility, descent obstruction, or update-vs-quotient
    mismatch. Generic semantics, generic predicate meaning, generic
    selection criteria, and generic agency are not admissible as
    \Pone{} projections in the absence of such a witness.
  \item \Ptwo{}: admissible only when the candidate carries a macro
    specification, a realization relation, specification-satisfaction
    data, and evidence of representability or non-representability.
    Macro-level predicates or objectives count only insofar as they are
    treated as content to be represented. Mere quotient equality,
    arbitrary relations without representability structure, and
    empirical observations without level and threshold data are not
    admissible as \Ptwo{} projections.
  \item \Pthree{}: admissible only when the candidate carries enriched
    route grammar, generated same-source, same-target route terms, and
    route-comparison data with a non-reducibility or coherence-mismatch
    certificate. Generic transitions, generic causal arrows, generic
    path records, and route undefinedness arising from \Pone{}
    obstruction are not admissible as \Pthree{} projections unless the
    route/coherence threshold is explicitly present and the route
    comparison is specifically of the \Pthree{} type.
  \item \Pfour{}: admissible only when the candidate carries a stage,
    order, transition, refinement, staged-dynamics, or regime-ordering
    record with laws beyond bare labels. Bare labels, arbitrary
    relations without order, transition, or refinement law, quotient
    packaging, and audit bookkeeping are not admissible as \Pfour{}
    projections.
  \item \Pfive{}: admissible only when the candidate carries a quotient
    relation, a packaging map, the identification of distinct
    micro-elements, a same-package or quotient-object record, and
    quotient-validity evidence. Generic structural paths, derivations,
    topologies or geometries merely because they have paths, and causal
    chains are not admissible as \Pfive{} projections unless quotient
    or packaging data is present.
  \item \Psix{}: admissible only when the candidate carries an audit
    record with event reference, provenance or dependency trace, replay
    or check structure, and bookkeeping or checkability structure.
    Passive logs alone, arbitrary metadata, and certificate-only
    content without audit/provenance/replay/check structure are not
    admissible as \Psix{} projections.
\end{itemize}
\end{proposition}

\begin{remark}[Three corrected alignments, held fixed]\label{rem:three-repaired-alignments}
Proposition~\ref{prop:alignment-rules} fixes three alignments that
are held throughout: generic transition data does not project to
\Pthree{}, generic path or derivation data does not project to
\Pfive{}, and generic semantic or interpretation data does not project
to \Pone{}. These are carried into every step of the classification
theorem below.
\end{remark}

\subsection{High-priority threat closures}\label{subsec:threat-closures}

Three classes of candidate seventh-role feature require particular
scrutiny: probability and uncertainty, residual causality in the
sense of the interventional and counterfactual
hierarchy~\citep{Pearl2009}, and residual agency in the planning
sense~\citep{Bratman1987}. Each closure is recorded as a proposition.
In every case the candidate is classified under the alignment rules
of Proposition~\ref{prop:alignment-rules} together with the residual
scheme of Definition~\ref{def:residual-scheme}; no candidate is
dismissed by definition. These closures remove the remaining
high-priority blockers before the general classification theorem.

\begin{proposition}[Closure of probability and uncertainty]\label{prop:probability-closure}
Every admissible probability or uncertainty feature of a description
\(D \in \FATCD{}\) classifies, within the covered scope, as one of
\begin{itemize}[topsep=2pt,itemsep=1pt]
  \item a \Ptwo{} projection, when the feature is a macro-specification
    or representability condition;
  \item a \Pfour{} projection, when the feature is a stochastic
    transition or regime structure carrying valid stage or transition
    data;
  \item a \Psix{} projection, when the feature is an audited
    uncertainty-bookkeeping record;
  \item an activation-threshold refinement, when the uncertainty data
    modulates when an already admissible role projection activates
    without supplying a new witness schema;
  \item an empirical-bridge annotation, when the feature expresses
    measurement uncertainty bridged to a substrate observable;
  \item or a composite of the foregoing projections and annotations when
    the admissible bookkeeping carries several of these aspects at once.
\end{itemize}
No residual probability or uncertainty feature remains as an unresolved
seventh-role threat within the covered scope.
\end{proposition}

\begin{proposition}[Closure of residual causality]\label{prop:causality-closure}
Every admissible residual-causality feature of a description
\(D \in \FATCD{}\) classifies, within the covered scope, as one of
\begin{itemize}[topsep=2pt,itemsep=1pt]
  \item a \Pfour{} projection, as the primary classification when the
    feature is a directed transition, refinement, or staged-dependence
    record;
  \item a \Pthree{} projection, only when the feature is specifically a
    route/coherence comparison between generated same-source,
    same-target routes;
  \item a \Psix{} projection, when the feature is a provenance or
    dependency audit;
  \item a composite of \Pfour{} and \Psix{}, when directed dependence
    and provenance audit occur together in the same admissible feature;
  \item an outside-domain structure, when the feature requires
    metaphysical causation not amenable to honest bookkeeping.
\end{itemize}
Classification as \Pfive{} is inadmissible in the absence of quotient
packaging. No residual-causality feature remains as an unresolved
seventh-role threat within the covered scope.
\end{proposition}

\begin{proposition}[Closure of residual agency]\label{prop:agency-closure}
Every admissible residual-agency feature of a description
\(D \in \FATCD{}\) classifies, within the covered scope, as one of
\begin{itemize}[topsep=2pt,itemsep=1pt]
  \item a composite of \Ptwo{}, \Pfour{}, and \Psix{}, when the feature
    decomposes into an objective or specification to be represented, a
    transition, regime, or staged selection event, and an audited
    decision-provenance record;
  \item an empirical-bridge annotation, when observed agency serves as
    substrate evidence for a bridged claim;
  \item an outside-domain structure, when irreducible intentional or
    teleological content is asserted beyond what admissible
    bookkeeping reduces to \Ptwo{}, \Pfour{}, and \Psix{}.
\end{itemize}
Classification as \Pone{} is inadmissible in the absence of an explicit
update-vs-quotient descent obstruction. No residual-agency feature remains
as an unresolved seventh-role threat within the covered scope.
\end{proposition}

\begin{remark}[Scope of the closures]\label{rem:threat-closures-scope}
Propositions~\ref{prop:probability-closure}--\ref{prop:agency-closure}
close the three threats within the covered scope; they do not make
scope-free claims. A feature that falls outside the covered scope is
classified as an outside-domain structure via
Definition~\ref{def:residual-scheme}, not absorbed into a primitive role
by default.
\end{remark}

\subsection{The upper-bound classification theorem}\label{subsec:classification-theorem}

The alignment rules of Proposition~\ref{prop:alignment-rules}, the
threat closures of
Propositions~\ref{prop:probability-closure}--\ref{prop:agency-closure},
and the residual scheme of Definition~\ref{def:residual-scheme}
combine to yield the upper-bound classification theorem: alignment
rules fix the admissible role meanings, threat closures rule out the
three misclassifications that otherwise block the argument, and the
theorem records the classification codomain for every covered
residual feature.

\begin{theorem}[Upper-bound classification]\label{thm:upper-bound-classification}
Let \(D \in \FATCD{}\) and let \(r\) be a residual feature of \(D\) in the
sense of Definition~\ref{def:residual-scheme}. Then within the covered
scope, \(r\) classifies as one of
\begin{enumerate}[label=(\roman*),topsep=2pt,itemsep=1pt]
  \item already absorbed into an admissible \Pone{}, \Ptwo{}, \Pthree{},
    \Pfour{}, \Pfive{}, or \Psix{} projection under the alignment rules of
    Proposition~\ref{prop:alignment-rules};
  \item a composite of such projections;
  \item a presentation variant of an already classified feature;
  \item a level shift of an already classified feature;
  \item an activation-threshold refinement;
  \item an instrument-visibility artifact;
  \item an empirical-bridge annotation;
  \item an outside-domain structure;
  \item already-covered bookkeeping or annotation data.
\end{enumerate}
Equivalently, within the covered scope, categories (x) and (xi) of
Definition~\ref{def:residual-scheme}, namely the unresolved-threat and
genuine-seventh-role categories, are empty: no covered candidate
seventh-role feature remains as an unresolved threat, and no covered
candidate seventh-role feature requires a genuine seventh irreducible
typed closure role.
\end{theorem}

\begin{proof}[Outline]
Let \(r\) be a residual feature of \(D\). If its raw feature form,
threshold data, witness data, level assignment, and bookkeeping align
with one of the fixed role meanings under
Proposition~\ref{prop:alignment-rules}, then \(r\) falls under category
(i). If several such aligned components are required, \(r\) falls under
the composite category (ii). Presentation variants and level shifts fall
under (iii) and (iv), activation-threshold refinements under (v),
instrument-visibility artifacts under (vi), empirical-bridge annotations
under (vii), and outside-domain structures under (viii). Features already
covered as bookkeeping or annotation data fall under (ix). The three
high-priority threat classes are handled decisively by
Propositions~\ref{prop:probability-closure}--\ref{prop:agency-closure},
so none of them remains in the unresolved-threat bucket. Therefore every
covered residual feature lands in categories (i)--(ix), and categories
(x) and (xi) are empty in the covered scope.
\end{proof}

\begin{remark}[Scope of the theorem]\label{rem:classification-scope}
Theorem~\ref{thm:upper-bound-classification} is scoped to \FATCD{}
and its covered scope. It does not claim that no seventh role can
arise outside \FATCD{}, that the alignment rules are exhaustive for
richer domains, or that the three threat closures settle analogous
questions beyond the covered scope. Nor does it prove exact-six on
its own: the decomposition/exhaustion theorem of
Section~\ref{sec:decomposition-exhaustion} remains necessary. Broader
questions are deferred to Section~\ref{sec:future-work}.
\end{remark}
